\documentclass[12pt,a4paper]{article}
% Drake_preamble.tex - LaTeX Preamble Template
% 作者: 謝慕揚, MD, PhD, FESC
% 
% 使用說明:
% 1. 表格請使用 longtable，不要有垂直分隔線 |
% 2. 表格內換行使用 \newline，不要使用 \\
% 3. 藥物名稱、檢驗、檢查請維持英文
% 4. Reference 格式請使用 \href 加上驗證過的 DOI

% Including essential packages for Chinese typesetting
\usepackage{xeCJK}
\usepackage{fontspec}
% Configuring Chinese font with Noto Serif CJK TC, fallback to SimSun
\IfFontExistsTF{Noto Serif CJK TC}{
  \setCJKmainfont{Noto Serif CJK TC}
  \setCJKsansfont{Noto Serif CJK TC}
  \setCJKmonofont{Noto Serif CJK TC}
}{\setCJKmainfont{SimSun}
  \setCJKsansfont{SimSun}
  \setCJKmonofont{SimSun}
}

% Including other necessary packages
\usepackage{geometry}
\usepackage{titlesec}
\usepackage{enumitem}
\usepackage{xcolor}
\usepackage{colortbl}
\usepackage{tcolorbox}
\usepackage{booktabs}
\usepackage{longtable}
\usepackage{array}
\usepackage{graphicx}
\usepackage{tikz}
\usepackage{fancyhdr}
\usepackage{multicol}
\usepackage{datetime2}
\usepackage{hyperref}
\usepackage{mdframed}
\usepackage{amsmath}
\usepackage{amssymb}
\usepackage{multirow}

\hypersetup{
    colorlinks=true,
    linkcolor=emergency_blue,
    citecolor=emergency_blue,
    urlcolor=emergency_blue,
    pdfborder={0 0 0}
}

% Defining page geometry
\geometry{
  top=2.5cm,
  bottom=2.5cm,
  left=2cm,
  right=2cm,
  headheight=25pt
}

% Adjusting topmargin to compensate for increased headheight
\addtolength{\topmargin}{-10pt}

% Defining colors
\definecolor{emergency_red}{RGB}{186,24,27}
\definecolor{emergency_blue}{RGB}{0,114,188}
\definecolor{emergency_gray}{RGB}{120,120,120}
\definecolor{highlight_yellow}{RGB}{255,250,205}

% Setting title formats
\titleformat{\section}[block]{\Large\bfseries\color{emergency_red}}{\thesection}{10pt}{}
\titleformat{\subsection}[block]{\large\bfseries\color{emergency_blue}}{\thesubsection}{8pt}{}
\titleformat{\subsubsection}[block]{\normalsize\bfseries\color{emergency_gray}}{\thesubsubsection}{6pt}{}

% Defining custom environment for important notes
\newmdenv[
    linecolor=emergency_red,
    backgroundcolor=highlight_yellow,
    linewidth=2pt,
    topline=true,
    bottomline=true,
    leftline=true,
    rightline=true,
    innertopmargin=10pt,
    innerbottommargin=10pt,
    innerleftmargin=10pt,
    innerrightmargin=10pt
]{importantbox}

% Defining note command with tcolorbox
\newcommand{\note}[1]{
    \par
    \noindent
    \begin{tcolorbox}[
        colback=emergency_blue!5!white,
        colframe=emergency_blue,
        title=\textbf{重點提醒},
        fonttitle=\bfseries,
        boxrule=0.5pt,
        arc=3pt,
        left=6pt,
        right=6pt,
        top=6pt,
        bottom=6pt
    ]
    #1
    \end{tcolorbox}
}

% Key findings box
\newcommand{\keyfinding}[1]{
    \par
    \noindent
    \begin{tcolorbox}[
        colback=yellow!10!white,
        colframe=orange!75!black,
        title=\textbf{核心發現摘要},
        fonttitle=\bfseries,
        boxrule=0.5pt,
        arc=3pt
    ]
    #1
    \end{tcolorbox}
}

% Defining custom date format for ISO 8601
\DTMnewdatestyle{iso}{
  \renewcommand{\DTMdisplaydate}[4]{\number##1-\number##2-\number##3}
  \renewcommand{\DTMDisplaydate}[4]{\number##1-\number##2-\number##3}
}
\DTMsetdatestyle{iso}
\newcommand{\mydate}{\DTMtoday}


% Custom header for this document
\pagestyle{fancy}
\fancyhf{}
\fancyhead[L]{\textcolor{emergency_red}{J Invasive Cardiol 教學講義}}
\fancyhead[R]{\textcolor{emergency_blue}{母子導管併發症與 OCT}}
\fancyfoot[C]{\textcolor{emergency_gray}{\thepage}}
\renewcommand{\headrulewidth}{1pt}
\renewcommand{\headrule}{\hbox to\headwidth{\color{emergency_red}\leaders\hrule height \headrulewidth\hfill}}

\begin{document}

%============================================================
% TITLE PAGE
%============================================================
\begin{center}
{\LARGE\bfseries\textcolor{emergency_red}{Journal of Invasive Cardiology - Clinical Images}}\\[0.5cm]
{\Large\textbf{Mother-and-Child Catheter-Induced Retrograde Dissection of the Left Main Coronary Artery During Optical Coherence Tomography Examination}}\\[0.3cm]
{\large 母子導管於光學同調斷層掃描檢查中引起左主幹冠狀動脈逆行性剝離}\\[0.5cm]
{\normalsize \href{https://pubmed.ncbi.nlm.nih.gov/29958179/}{J Invasive Cardiol. 2018;30(7):E57-E58.}}\\[0.3cm]
{\normalsize 整理：謝慕揚 MD, PhD, FESC \quad 日期：\mydate}
\end{center}

\vspace{0.5cm}

%============================================================
\section{病例報告}
%============================================================

\subsection{病人基本資料}

\begin{itemize}
    \item 54 歲男性
    \item 診斷：Non-ST elevation myocardial infarction (NSTEMI)
\end{itemize}

\subsection{冠狀動脈攝影發現}

冠狀動脈攝影顯示雙血管疾病 (two-vessel disease)：

\begin{longtable}{p{5cm}p{9cm}}
\toprule
\textbf{血管} & \textbf{病灶} \\
\midrule
\endfirsthead
\toprule
\textbf{血管} & \textbf{病灶} \\
\midrule
\endhead
Left anterior descending (LAD) & Mid LAD 嚴重狹窄 \\
\midrule
Left circumflex (LCX) & Proximal 及 distal LCX 嚴重狹窄 \\
\bottomrule
\end{longtable}

%============================================================
\section{介入治療過程}
%============================================================

\subsection{初始介入策略}

\begin{itemize}
    \item 選擇先處理 LCX 病灶
    \item 在 optical coherence tomography (OCT) 影像導引下進行經皮冠狀動脈介入術 (PCI)
    \item 植入 3 × 40 mm drug-eluting stent (DES) 於 LCX
\end{itemize}

\subsection{支架植入併發症}

\begin{importantbox}
\textbf{問題 1：支架位置偏移}

由於病人呼吸動作影響，支架著陸位置比預期更近端 (more proximally than intended)。

OCT 顯示近端支架嚴重貼壁不良 (gross malapposition)。
\end{importantbox}

\subsection{使用 Mother-and-Child 導管系統}

\subsubsection{使用原因}

為了使 4 mm non-compliant balloon 能順利通過貼壁不良的支架進行後擴張：

\begin{itemize}
    \item 將 6 Fr GuideLiner「child」支撐導管插入現有「mother」6 Fr guiding catheter
    \item 目的：提供額外的同軸支撐力 (coaxial back-up support)
\end{itemize}

\subsubsection{OCT 影像問題}

支架後擴張後的 OCT 影像品質不佳，原因：

\begin{itemize}
    \item 顯影劑優先流入較大的 LAD 血管
    \item 解決方案：將 GuideLiner 軟尖端推進至 proximal LCX 以進行選擇性顯影劑注射
\end{itemize}

\note{
\textbf{OCT 對比劑注射慣例}：該單位常規使用極短暫的手動注射 10-12 mL 顯影劑進行 OCT 影像擷取。
}

%============================================================
\section{嚴重併發症：醫源性冠狀動脈剝離}
%============================================================

\begin{importantbox}
\textbf{災難性併發症}

OCT 影像擷取後的血管攝影顯示：

\textbf{Hydraulic spiral dissection}（水力螺旋剝離）
\begin{itemize}
    \item 起始於 proximal LCX
    \item 逆行性延伸至 proximal LAD
    \item 進一步延伸至 distal left main (LM) stem
\end{itemize}
\end{importantbox}

\subsection{發生機轉分析}

\begin{enumerate}
    \item GuideLiner 的軟尖端未與迂曲的 LCX 血管同軸
    \item 導管尖端可能抵住血管壁或斑塊
    \item 儘管注射顯影劑前無壓力衰減 (pressure damping) 警訊
    \item 短暫且有力的手動顯影劑注射產生的局部水力壓力 (focal hydraulic pressure)
    \item 在脆弱斑塊處造成裂縫並沿損傷點追蹤，導致冠狀動脈剝離
\end{enumerate}

%============================================================
\section{併發症處理}
%============================================================

\subsection{緊急處置步驟}

\begin{longtable}{p{1cm}p{13cm}}
\toprule
\textbf{步驟} & \textbf{處置內容} \\
\midrule
\endfirsthead
\toprule
\textbf{步驟} & \textbf{處置內容} \\
\midrule
\endhead
1 & 立即在 LCX 剝離處充氣 4.0 mm semi-compliant (SC) balloon，維持 1 分鐘 \\
\midrule
2 & 球囊阻塞 LCX 時進行血管攝影，確認 LAD 通暢且無顯影劑滲漏 (contrast staining) \\
\midrule
3 & 從 LM 至 LCX 植入 4 × 24 mm DES，與初始 LCX 支架重疊 \\
\midrule
4 & 最終血管攝影結果滿意 \\
\bottomrule
\end{longtable}

\subsection{病人結果}

\begin{itemize}
    \item 整個手術過程中病人維持穩定
    \item 成功救治醫源性左主幹剝離
\end{itemize}

%============================================================
\section{Mother-and-Child 導管系統介紹}
%============================================================

\subsection{GuideLiner 支撐導管}

\begin{longtable}{p{4cm}p{10cm}}
\toprule
\textbf{特性} & \textbf{說明} \\
\midrule
\endfirsthead
\toprule
\textbf{特性} & \textbf{說明} \\
\midrule
\endhead
製造商 & Vascular Solutions \\
\midrule
設計 & 「Child」導管通過標準「mother」guiding catheter 輸送 \\
\midrule
功能 & 允許對目標血管進行深度插管 (deep intubation) \\
\bottomrule
\end{longtable}

\subsection{臨床應用}

\begin{enumerate}
    \item \textbf{增強支撐力}：提供同軸支撐以利支架輸送
    \item \textbf{困難病灶處理}：
    \begin{itemize}
        \item 嚴重鈣化病灶
        \item 迂曲血管
    \end{itemize}
    \item \textbf{選擇性顯影劑注射}：改善 OCT 等腔內影像品質
\end{enumerate}

%============================================================
\section{教學重點}
%============================================================

\subsection{OCT 影像擷取時的風險因素}

\begin{importantbox}
\textbf{使用 GuideLiner 進行 OCT 時的注意事項}

\begin{enumerate}
    \item \textbf{導管位置}：
    \begin{itemize}
        \item GuideLiner 尖端必須與血管同軸
        \item 在迂曲血管中尤需特別注意
    \end{itemize}

    \item \textbf{壓力監測限制}：
    \begin{itemize}
        \item 無壓力衰減不代表導管位置安全
        \item 導管可能仍抵住血管壁或斑塊
    \end{itemize}

    \item \textbf{顯影劑注射風險}：
    \begin{itemize}
        \item 短暫有力的手動注射可產生高局部水力壓力
        \item 可在脆弱斑塊處造成裂縫並引發剝離
    \end{itemize}
\end{enumerate}
\end{importantbox}

\subsection{預防醫源性剝離的策略}

\begin{enumerate}
    \item 在迂曲血管中定位 GuideLiner 需格外小心
    \item 確保導管尖端與血管同軸
    \item 考慮使用較緩和的顯影劑注射方式
    \item 注射前仔細評估導管位置
    \item 若影像品質不佳，考慮替代策略而非強行深度插管
\end{enumerate}

\subsection{Catheter-induced coronary dissection 的處理原則}

\begin{longtable}{p{4cm}p{10cm}}
\toprule
\textbf{處理策略} & \textbf{說明} \\
\midrule
\endfirsthead
\toprule
\textbf{處理策略} & \textbf{說明} \\
\midrule
\endhead
立即球囊充氣 & 在剝離處充氣球囊以壓迫剝離瓣膜，阻止剝離延伸 \\
\midrule
評估血流 & 確認剝離未影響其他重要血管的血流 \\
\midrule
支架植入 & 從正常血管段至剝離區域植入支架以封閉剝離 \\
\midrule
持續監測 & 確保血流動力學穩定及最終結果滿意 \\
\bottomrule
\end{longtable}

%============================================================
\section{結論}
%============================================================

\note{
\textbf{Take-home Message}

Mother-and-child 導管系統在迂曲血管中進行 OCT 影像擷取時，需特別注意 GuideLiner 的定位，以避免醫源性冠狀動脈剝離。即使無壓力衰減警訊，導管尖端的非同軸位置加上有力的顯影劑注射，仍可在脆弱斑塊處造成剝離。
}

%============================================================
\section{參考文獻}
%============================================================

\begin{enumerate}
    \item Chi WK, Yan BP, Chan CY. Mother-and-Child Catheter-Induced Retrograde Dissection of the Left Main Coronary Artery During Optical Coherence Tomography Examination. \href{https://pubmed.ncbi.nlm.nih.gov/29958179/}{\textit{J Invasive Cardiol}. 2018;30(7):E57-E58.}

    \item Parise H, Maehara A, Stone GW, et al. Meta-analysis of randomized studies comparing intravascular ultrasound versus angiographic guidance of percutaneous coronary intervention in pre-drug-eluting stent era. \textit{Am J Cardiol}. 2011;107:374-382.

    \item Ali ZA, Maehara A, Genereux P, et al. Optical coherence tomography compared with intravascular ultrasound and with angiography to guide coronary stent implantation (ILUMIEN III: OPTIMIZE PCI): a randomised controlled trial. \textit{Lancet}. 2016;388:2618-2628.

    \item Vo MN, Mintz GS, McPherson JA. GuideLiner catheter facilitated percutaneous coronary intervention. \textit{J Invasive Cardiol}. 2013;25:E191-E197.
\end{enumerate}

\end{document}
